\documentclass[a4paper,12pt]{article} 
\usepackage[utf8]{inputenc} % Acentos válidos sin problemas
\usepackage[spanish]{babel} % Idioma
%-----------------------------------INICIO DE PACKETES------------------/
%----------------------------------------------------------------------/|
\usepackage{amsmath}   % Matemáticas: Comandos extras(cajas ecuaciones) |
\usepackage{amssymb}   % Matemáticas: Símbolos matemáticos              |
\usepackage{datetime}  % Agregar fechas                                 |
\usepackage{graphicx}  % Insertar Imágenes                              |
%\usepackage{biblatex} % Bibliografía                                   |
\usepackage{multicol}  % Creación de tablas                             |
\usepackage{longtable} % Tablas más largas                              |
\usepackage{xcolor}    % Permite cambiar colores del texto              |
\usepackage{tcolorbox} % Cajas de color                                 |
\usepackage{setspace}  % Usar espacios                                  |
\usepackage{fancyhdr}  % Para agregar encabezado y pie de página        |
\usepackage{float}     % Flotantes                                      |
\usepackage{soul}      % "Efectos" en palabras                          |
\usepackage{hyperref}  % Para usar hipervínculos                        |
\usepackage{caption}   % Utilizar las referencias                       |
\usepackage{subcaption} % Poder usar subfiguras                         |
\usepackage{multirow}  % Nos permite modificar tablas                   |
\usepackage{array}     % Permite utilizar los valores para multicolumn  |
\usepackage{booktabs}  % Permite modificar tablas                       |
\usepackage{diagbox}   % Diagonales para las tablas                     |
\usepackage{colortbl}  % Color para tablas                              |
\usepackage{listings}  % Escribir código                                |
\usepackage{mathtools} % SIMBOLO :=                                     |
\usepackage{enumitem}  % Modificar items de Listas                      |
\usepackage{tikz}      %                                                |
\usepackage{lipsum}    % for auto generating text                       |
\usepackage{afterpage} % for "\afterpage"s                              |
\usepackage{pagecolor} % With option pagecolor={somecolor or none}|     | 
%\usepackage{glossaries} %                                              |        
%----------------------------------------------------------------------\|
%-----------------------------------FIN--- DE PACKETES------------------\
%\usepackage{newtxtext,newtxmath}
%\usepackage{bm}
%--------------------------------/
%-------------------------------/
\usepackage[                 %   |
  headheight=15pt,  %            |
  letterpaper,  % Tipo de pag.   |
  left =1.5cm,  %  < 1 >         |
  right =1.5cm, %  < 1 >         | MARGENES DE LA PAGINA
  top =2cm,     %  < 1.5 >       |
  bottom =1.5cm %  < 1.5 >       |
]{geometry}     %                |
%-------------------------------\
%--------------------------------\

%----------------------------------------------------------------------/
%-------------------Encabezado y Pie de Pagina-----------------------/ |
%--------------------------------------------------------------------\ |
\usepackage{fancyhdr} %                                                |
\usepackage{lastpage} %                                                |
%\fancyhf{}           %                                                |
\pagestyle{fancy}     %                                                |
\fancyhead[L]{{$2024-1$}} %                                              |
\fancyhead[C]{}                      
\fancyhead[R]{}

\fancyfoot[L]{} 
\fancyfoot[C]{\thepage}
\fancyfoot[R]{} 

\renewcommand{\headrulewidth}{1.5pt} %                                 |
\renewcommand{\footrulewidth}{1.5pt} %                                 | 
%--------------------------------------------------------------------\ |
%-------------------Encabezado y Pie de Pagina-----------------------/ |
%----------------------------------------------------------------------/


\begin{document}%----------------------INICIO DOCUMENTO------------|
%------------------------------------------------------------------|

\begin{titlepage}
    \center% Center everything on the page
    % Defines a new command for the horizontal lines, change thickness here
    \newcommand{\HRule}{\rule{\linewidth}{0.5mm}} 
    
    %---------------------
    %	LOGO
    %---------------------
    \includegraphics[width=4.5cm]{IMA/Ciencias.png} \\ 
    
    %----------------------------
    %	Encabezado
    %----------------------------
    \quad \\[1.5cm]
    %\textsc{\LARGE MSc Thesis}\\[1.5cm] % Name of your university/college
    \textsc{\Large Facultad de Ciencias}\\[0.5cm] % Major heading such as course name
    \textsc{\Large Computación Distribuida}\\[0.5cm] % Major heading such as course name
    %\textsc{\large Del Aula a la START UP}\\[0.5cm] % Minor heading such as course title
    %-------------
    %	TITULOS
    %-------------
    \makeatletter
        \HRule\\ [0.4cm]
            { \huge \bfseries PRACTICA 02}\\
        \HRule\\ [0.4cm]
    %-------------------
    %	PROFESORES
    %-------------------
    %\vspace{10mm}
    \large\textbf{Semestre $2024-1$}\\
    
    \vspace{10mm}
    \begin{minipage}{0.7\textwidth}
        \begin{flushleft} \large
            \emph{Profesor:}\\
                Luis Germán Pérez Hernández\\
                \vspace{3mm}
                \emph{Ayudantes:}\\
                Daniel Michel Tavera\\
                Yael Antonio Calzada Martín\\
        \end{flushleft}
    \end{minipage}
    %~
    \\
    %-------------------
    %	AUTOR
    %-------------------
    \vspace{10mm}
    \begin{minipage}{0.4\textwidth}
        \begin{flushright} \large
        \centering
            \emph{Alumnos} \\ % \emph{Alumnos:} \\
            \vspace{3mm}
            Marco Silva Huerta\\
            316205326\\
            \vspace{4mm}
            Edgar Montiel Ledesma\\
            317317794\\
            \vspace{4mm}
            Carlos Cortés\\
            420004846
            %\vspace{4mm}
            % Autor 2\\
        \end{flushright}
    \end{minipage}\\[1.2cm]
    \makeatother
    %---------------
    %	F E C H A
    %---------------
    {\large 18 de Octubre de 2023}\\[2cm]
    \vfill 
    \end{titlepage}
\newpage

\begin{center}
    {\huge Algoritmo abusón (bully)}
\end{center}

%%%%%%%%%%%%%%%%%%%%%%%%%%%%%%%%%%%%%%%
%%%%%%%%%%%%%%%%%%%%%%%%%%%%%%%%%%%%%%
%%%%%%%%%%%%%%%%%%%%%%%%%%%%%%%%%%%%
\section*{Ejecución del programa}
%%%%%%%%%%%%%%%%%%%%%%%%%%%%%%%%%%%%
%%%%%%%%%%%%%%%%%%%%%%%%%%%%%%%%%%%%%%
%%%%%%%%%%%%%%%%%%%%%%%%%%%%%%%%%%%%%%%

\subsection*{Compilar}

\begin{small}
    \begin{verbatim}
mpicc Practica03_EdgarMontiel_CarlosCortes_MarcoSilva.c -o Practica03_EdgarMontiel_CarlosCortes_MarcoSilva
    \end{verbatim}            
\end{small}
\subsection*{Ejecutar}
\begin{verbatim}
    ./Practica03_EdgarMontiel_CarlosCortes_MarcoSilva
\end{verbatim}


%%%%%%%%%%%%%%%%%%%%%%%%%%%%%%%%%%%%%%%
%%%%%%%%%%%%%%%%%%%%%%%%%%%%%%%%%%%%%%
%%%%%%%%%%%%%%%%%%%%%%%%%%%%%%%%%%%%
\section*{Elección distribuida}
%%%%%%%%%%%%%%%%%%%%%%%%%%%%%%%%%%%%
%%%%%%%%%%%%%%%%%%%%%%%%%%%%%%%%%%%%%%
%%%%%%%%%%%%%%%%%%%%%%%%%%%%%%%%%%%%%%%

\subsection*{Consideraciones}

\begin{itemize}
    \item Permite la caída de procesos durante la elección
     \begin{itemize}
        \item Utiliza \textit{timeouts} para detectar fallos de procesos
     \end{itemize}
    \item Supone comunicación fiable
    \item Cada proceso conoce qué procesos tienen identificadores mayores y puede comunicarse con ellos
\end{itemize}

\subsection*{Funcionamiento resumido}

\begin{enumerate}
    \item El convocante envía \textit{mensajes elección} a los procesos de $id$ mayor
    \item Si ninguno le responde, \textit{multidifunde} que es el nuevo coordinador
    \item Si alguno le responde, el convocante inicial queda en espera, y los procesos que responden
         inician un nuevo proceso de elección como convocantes (vuelta al paso 1)
\end{enumerate}

\subsection*{Tipos de mensaje}

\begin{itemize}
    \item \textbf{Elección:} anuncia un proceso de elección
    \item \textbf{Respuesta:} respuesta a un mensaje de elección
    \item \textbf{Coordinador:} anuncia la identidad del proceso elegido
\end{itemize}


%%%%%%%%%%%%%%%%%%%%%%%%%%%%%%%%%%%%%%%
%%%%%%%%%%%%%%%%%%%%%%%%%%%%%%%%%%%%%%
%%%%%%%%%%%%%%%%%%%%%%%%%%%%%%%%%%%%
\section*{Funcionamiento}
%%%%%%%%%%%%%%%%%%%%%%%%%%%%%%%%%%%%
%%%%%%%%%%%%%%%%%%%%%%%%%%%%%%%%%%%%%%
%%%%%%%%%%%%%%%%%%%%%%%%%%%%%%%%%%%%%%%



%%%%%%%%%%%%%%%%%%%%%%%%%%%%%%%%%%%%%%%
%%%%%%%%%%%%%%%%%%%%%%%%%%%%%%%%%%%%%%
%%%%%%%%%%%%%%%%%%%%%%%%%%%%%%%%%%%%
\section*{Pseudocódigo del Algoritmo}
%%%%%%%%%%%%%%%%%%%%%%%%%%%%%%%%%%%%
%%%%%%%%%%%%%%%%%%%%%%%%%%%%%%%%%%%%%%
%%%%%%%%%%%%%%%%%%%%%%%%%%%%%%%%%%%%%%%

\begin{verbatim}
    Algoritmo Abusivo

    Variables Globales:
    - Mi_Id: identificador único del proceso
    - Coordinador_Actual: identificador del coordinador actual
    - Esperando_Coordinador: booleano que indica si el proceso 
       está esperando un mensaje de coordinador
    
    Inicio del Algoritmo:

    Declaración de una estructura de Nodo 
        - Define una estructura llamada Nodo que almacena información sobre un nodo, incluyendo su ID, su estado de estar vivo o muerto, el ID de su líder y un indicador de si está esperando una respuesta.
    
    Cuando un proceso detecta que el coordinador ha fallado debido a un timeout:
        - Iniciar la elección:
        
        Si Mi_Id es el identificador más alto no fallido:
            - Mi_Id es el nuevo coordinador
            - Enviar un mensaje de coordinador a todos los procesos con 
               identificadores más bajos (proceso abusón)
    
        Si Mi_Id no es el identificador más alto no fallido:
            - Enviar un mensaje elección a todos los procesos con 
               identificadores más altos
            - Esperar un mensaje de respuesta durante un tiempo T
            - Si no se recibe ningún mensaje de respuesta después de un tiempo T:
                - Regresar al paso de Inicio del Algoritmo
            - Si se recibe un mensaje de respuesta:
                - Esperar un mensaje de coordinador
                - Si se recibe un mensaje de coordinador:
                    - Coordinador_Actual = Id contenido en el mensaje de 
                       coordinador
                - Si no se recibe un mensaje de coordinador:
                    - Iniciar una nueva elección

    Función simulaTimeout
        - Esta función simula un timeout con una probabilidad del 30% de fallar. Devuelve true si se produce un timeout y false en caso contrario.
    
    Si un proceso se recupera o se lanza un proceso sustituto con el mismo Mi_Id:
        - Iniciar una nueva elección, independientemente de si el coordinador 
           actual está funcionando

    Función main:
    - En la función principal, se inicializa la semilla aleatoria y se inicia MPI.
    - Se crea un arreglo de estructuras Nodo llamado nodos, y se inicializan los nodos con sus IDs, estado de estar vivo y líderes.
    - Se define un valor para myId, que representa el ID del nodo actual.
    - Se imprime un mensaje para indicar el inicio de la elección.
    - Se llama a la función iniciarElección para el nodo actual.
    - Luego, se imprime el resultado de la elección para cada nodo en el sistema.
    - Finalmente, se llama a MPI_Finalize para finalizar la comunicación MPI.
\end{verbatim}

%-----------------------------------------------------------------|
\end{document}%---------------------- FIN DOCUMENTO---------------|
%-----------------------------------------------------------------|